\section{Perspectiva unificada: el ciclo cerrado de automejora}

Esta clase puede resumirse en un ciclo de cuatro etapas:

\begin{enumerate}
    \item \textbf{Generate}: el modelo o varios Agents producen acciones, soluciones o planes candidatos;
    \item \textbf{Verify}: pruebas, un reward model, reglas o señales del entorno evalúan a los candidatos;
    \item \textbf{Select / Act}: se selecciona un candidato y se ejecuta en el entorno;
    \item \textbf{Update}: la retroalimentación se usa para corregir la trayectoria actual, actualizar la memoria o entrenar los parámetros.
\end{enumerate}

\begin{importantbox}{Mapa conceptual de todo el curso}
El test-time scaling mejora principalmente Generate y Select; la robust verification mejora Verify; el tool use y el planning mejoran Act; el RL y los synthetic data reescriben la retroalimentación en los parámetros; la memory y el continual learning conservan la experiencia entre tareas. Lo que llamamos self-improving agent es precisamente lograr que esta cadena pueda repetirse de manera estable y medible.
\end{importantbox}

\subsection{Lista de verificación al diseñar un sistema de automejora}

\begin{itemize}
    \item ¿De dónde viene la diversidad de los candidatos? ¿No se está simplemente repitiendo el mismo error?
    \item ¿El verifier es más confiable que el generator, o comparten el mismo punto ciego?
    \item ¿La retroalimentación es inmediata, está diferida, o solo puede aproximarse con una métrica proxy?
    \item ¿Qué acciones son reversibles y cuáles requieren permisos o confirmación humana?
    \item ¿El beneficio marginal del cómputo adicional compensa el costo y la latencia?
    \item ¿La actualización podría provocar reward hacking, olvido o deriva de las capacidades?
\end{itemize}

\subsection{Resumen de la sección}

Al separar el modelo, el verificador, las herramientas y el mecanismo de actualización, es posible ubicar con mayor claridad los cuellos de botella del sistema. Muchos fallos que parecen deberse a que "el modelo no es suficientemente inteligente" en realidad provienen del verificador, la gestión de estado, la interfaz de herramientas o la asignación del presupuesto, y no necesariamente requieren volver a entrenar un modelo base más grande.

